\chapter{量子力学初步}
\problem{3.1}{波长为 $1\,\angstrom$ 的 X 光光子的动量和能量各为多少？}
\answerarea
\problem{3.2}{经过 10000 伏特电势差加速的电子束的德布罗意波波长 $\lambda$ 为多少？用上述电压加速的质子束，其德布罗意波波长是多少 $\angstrom$？}
\answerarea\clearpage
\problem{3.3}{电子被加速后的速度很大，必须考虑相对论修正。因而原来
\[
\lambda=\frac{12.25}{\sqrt V}\;(\angstrom)
\]
的电子德布罗意波波长与加速电压的关系式应改为
\[
\lambda=\frac{12.25}{\sqrt V}\left(1-0.489\times10^{-6}V\right)\,\angstrom.
\]
其中 $V$ 为以伏特为单位的电子加速电压。试证明之。}
\answerarea
\problem{3.4}{试证明氢原子稳定轨道上正好能容纳整数个电子的德布罗意波波长（习题图 3.1）。上述结果不但适用于圆轨道，同样适用于椭圆轨道，试证明之。
\fig{3_1}{3.1}}
\answerarea\clearpage
\problem{3.5}{带电粒子在威耳孙云室（一种径迹探测器）中的轨迹是一串小雾滴，雾滴的线度约为 1 微米。当观察能量为 1000 电子伏特的电子径迹时其动量与经典力学动量的相对偏差不小于多少？}
\answerarea
\problem{3.6}{试证明自由运动的粒子（势能 $V\equiv0$）的能量可以有连续的值。}
\answerarea\clearpage
\problem{3.7}{粒子位于一维对称势场中，势场形式如下图，即
\[
V(x)=\begin{cases}0,&0<x<L,\\ V_0,&x<0\text{ 或 }x>L.\end{cases}
\]
\fig{3_2}{3.2}
\begin{enumerate}
\item 试推导粒子在 $E<V_0$ 情况下其总能量 $E$ 满足的关系式。
\item 试利用上述关系式，以图解法证明，粒子的能量只能是一些不连续的值。
\end{enumerate}}
\answerarea
\problem{3.8}{有一粒子，其质量为 $m$，在一个三维势箱中运动。势箱的长、宽、高分别为 $a,b,c$。在势箱外，势能 $V=\infty$；在势箱内，$V\equiv0$。试算出粒子可能具有的能量。}
\answerarea\clearpage
